\section{Distinguishability Before Sameness}

The first grammatical act is not naming but telling apart. A name
already asks too much. It asks that a sound, mark, or position be
allowed to return to itself; it asks that a reader recognize the return;
it asks that the record preserve what the first utterance placed in it.
At the foundation those returns have not been justified. The grammar has
not yet been granted an identity relation, a storehouse of objects, or a
rule that says two presentations are the same presentation. It can only
permit a discrimination.

This is why distinguishability is primitive. It is weaker than equality
and stronger than silence. It does not say that one side has an essence
and the other side has another essence. It says only that the two sides
cannot be read as one under the current act. The act is local: it holds
for this reading, under this available contrast, with this much structure
carried forward. It forbids the reader from smuggling in a universal
sameness before the record has earned the right to use one.

The difference is therefore not a property added to already completed
objects. It is the first way anything can enter the ledger. Before the
ledger distinguishes, there is no admissible place to write an object at
all. A mark that cannot be told apart from absence is not yet a mark. A
signal that cannot be separated from background is not yet a signal. A
state that cannot be discriminated from its complement is not yet a
state. Distinguishability is the permission by which a world becomes
recordable.

This reverses the usual order of comfort. The comfortable order says:
first there are things, then the things may be compared. The grammatical
order says: first a comparison is permitted, and only afterward may
things be stabilized as the poles of that comparison. An instrument
illustrates the point whenever it registers a threshold. The threshold
does not need to know the full nature of what crossed it. It needs only a
recoverable difference between crossing and not crossing. Later readings
may identify a particle, a voltage, a pressure, a word, or a color. The
first reading merely separates.

Relational examples make the same point. Precedence, inclusion, and
opposition are readable before they become algebraic operations. A child
can sort two blocks before naming the class to which either belongs. A
scale can separate heavier from lighter before the unit of mass is
calibrated. A sensor can say that a new record differs from the old
record before any theory has decided what substance, force, or law the
difference represents. The grammar permits the relation before it
identifies the relata as finished objects.

Nothing in this permission yet licenses truth. To read a difference is
not to say which side ought to be affirmed. The grammar has only made a
space in which affirmation could later have a target. It has also made a
space in which denial could later have a target. Because both spaces
arise from the same primitive permission, neither can claim priority at
the foundation. This will matter at the collapse. If affirmation and
denial later ask the bare grammar to distinguish them from the inside,
the grammar can only use what has already been granted. It has been
granted difference, not self-standing truth.

The first permission is therefore austere but fertile. It permits a
separation; it forbids unearned identity; it carries enough structure for
a second act. Counting is that second act. But the count will not be a
smooth ascent from obvious units. It will inherit the tension in the
first distinction: the grammar may read across a sameness it has
manufactured for the occasion, or it may read the fracture that keeps the
sides apart. Wholes and parts are already waiting inside the first
difference.
